\section{Experiments}
\label{sec:exp}

We evaluate KG-SCoRE on four questions: (Q1) Does LLM-based KG verification achieve competitive recall \emph{and} faithfulness compared to baselines under fair evaluation? (Q2) What is the contribution of each module, and does ComplEx scoring help or hurt? (Q3) Is the faithfulness metric reliable? (Q4) Does LLM verification actually clean the KG?

\subsection{Setup}
\label{sec:exp:setup}

\paragraph{Knowledge graph.}
We use a SemMedDB-style~\cite{kilicoglu2012semmeddb} biomedical KG with 9{,}188 triples over 523 entities (diseases, chemicals, genes, phenotypes) and 32 relation types. The KG has a power-law degree distribution with typed entity constraints; we inject $\sim$35\% noisy type-violating triples to simulate SemMedDB extraction errors~\cite{kilicoglu2020broadcoverage}. Since we know which triples are noise (score $<0.5$), we can directly evaluate KG cleaning accuracy (Q4).

\paragraph{Gold standard.}
We evaluate on the Swanson--Cameron LBD benchmark~\cite{hristovski2015constructing,bhasuran2023literature}: 20 disease--chemical pairs. The first 10 are classic Swanson pairs; the remaining 10 are extension pairs. A method succeeds if the gold chemical appears in its top-$k$ hypotheses.

\paragraph{Baselines.}
(1) \textbf{Vanilla LLM}: no retrieval, no KG. (2) \textbf{BM25-RAG}: BM25 snippet retrieval + LLM. (3) \textbf{ToG}~\cite{sun2023thinkongraph}: LLM-guided KG exploration. (4) \textbf{RoG}~\cite{luo2024graphconstrained}: path-constrained retrieval + LLM generation. All use GPT-4o-mini (temperature 0.2).

\paragraph{Metrics.}
Re-discovery \textbf{Hit@$k$}, \textbf{MRR}, \textbf{StrictFaith@$k$} (Eq.~\ref{eq:faith_strict}), \textbf{BroadFaith@$k$} (Eq.~\ref{eq:faith_broad}), and \textbf{Path Coverage} (fraction of hypotheses with extractable KG paths). For fair faithfulness evaluation, all methods receive post-hoc path extraction: for methods that do not emit paths, we extract paths by matching hypothesis entities against the retrieved subgraph. We run 3 seeds and report McNemar tests for binary Hit@10.

\paragraph{ComplEx training.}
ComplEx ($d{=}200$, 100 epochs, lr 0.1) via PyKEEN on a single RTX 5090 ($\sim$4 min). Held-out MRR is 0.013---above chance (1/523 $\approx$ 0.002) but modest, reflecting the KG's high noise ratio. This low MRR is itself a finding: ComplEx scores are unreliable for triple filtering on noisy KGs.

\subsection{Main Results (Q1)}
\label{sec:exp:main}

Table~\ref{tab:main} reports LBD re-discovery and fair faithfulness on 20 pairs.

\begin{table}[tb]
  \centering
  \caption{LBD re-discovery and fair faithfulness on 9{,}188-triple KG (20 pairs). Best in \textbf{bold}.}
  \label{tab:main}
  \small
  \begin{tabular}{lccccc}
    \toprule
    Method & Hit@10 & MRR & StrictFaith@10 & BroadFaith@10 & Coverage \\
    \midrule
    Vanilla LLM & 0.25 & 0.08 & 0.00 & 0.03 & 0.11 \\
    BM25-RAG & 0.00 & 0.00 & 0.00 & 0.17 & 0.57 \\
    ToG & 0.10 & 0.03 & 0.05 & 0.94 & 1.00 \\
    RoG & \textbf{0.40} & \textbf{0.18} & 0.04 & \textbf{0.99} & \textbf{1.00} \\
    \textbf{KG-SCoRE} & 0.25 & 0.15 & \textbf{0.10} & \textbf{0.99} & \textbf{1.00} \\
    \bottomrule
  \end{tabular}
\end{table}

\paragraph{Finding 1: Fair faithfulness reveals real differences.}
Unlike our earlier evaluation (which assigned 0 to all baselines), the fair evaluation reveals that ToG achieves BroadFaith@10 of 0.94 and Coverage of 1.00---it does explore and output KG paths. However, ToG's \emph{strict} faithfulness is only 0.05, meaning most hypotheses are only partially supported by their paths. RoG shows a similar pattern: BroadFaith@10 of 0.99 and Coverage of 1.00, but StrictFaith@10 of only 0.04. KG-SCoRE achieves StrictFaith@10 of 0.10 (2.5$\times$ RoG's and 2$\times$ ToG's) while matching BroadFaith@10 of 0.99. Vanilla LLM has near-zero coverage (0.11), confirming it generates without KG grounding. BM25-RAG has moderate coverage (0.57) but near-zero faith, indicating retrieved text snippets rarely align with KG paths. The pattern across all KG-grounded methods is consistent: broad alignment with paths is easy, but strict entailment is rare---KG-SCoRE's verification step doubles the strict faithfulness rate.

\paragraph{Finding 2: No significant recall difference from RoG.}
On 20 pairs, RoG achieves the highest Hit@10 (0.40), with KG-SCoRE and Vanilla LLM tied at 0.25. On the 10-pair multi-seed subset (Table~\ref{tab:multiseed}), KG-SCoRE achieves 0.47$\pm$0.05 vs.\ RoG's 0.57$\pm$0.05. McNemar test shows no significant difference ($p{>}0.20$), indicating failure to detect a recall gap---not proof of equivalence.

\paragraph{Statistical significance.}
McNemar tests on 10-pair multi-seed ($n{=}30$) confirm KG-SCoRE significantly outperforms Vanilla LLM ($p{<}0.01$), BM25-RAG ($p{<}0.05$), and ToG ($p{<}0.05$), but not RoG ($p{>}0.20$).

\begin{table}[tb]
  \centering
  \caption{Multi-seed Hit@10 ($n=30$, 3 seeds $\times$ 10 pairs). Seeds vary LLM sampling randomness at fixed temperature 0.2.}
  \label{tab:multiseed}
  \small
  \begin{tabular}{lcc}
    \toprule
    Method & Mean $\pm$ Std & Per-seed \\
    \midrule
    Vanilla LLM & 0.10 $\pm$ 0.14 & [0.00, 0.30, 0.00] \\
    BM25-RAG & 0.17 $\pm$ 0.05 & [0.20, 0.20, 0.10] \\
    ToG & 0.20 $\pm$ 0.00 & [0.20, 0.20, 0.20] \\
    \textbf{RoG} & \textbf{0.57 $\pm$ 0.05} & [0.60, 0.60, 0.50] \\
    \textbf{KG-SCoRE} & \textbf{0.47 $\pm$ 0.05} & [0.50, 0.40, 0.50] \\
    \bottomrule
  \end{tabular}
\end{table}

\subsection{Ablation: LLM Verification vs.\ Learned Scoring (Q2)}
\label{sec:exp:ablation}

Table~\ref{tab:ablation} isolates each module's contribution. The primary configuration is \texttt{+llm+faith} (LLM verification + faithfulness, no ComplEx).

\begin{table}[tb]
  \centering
  \caption{Ablation on 9{,}188-triple KG (10 pairs). \texttt{+llm+faith} is the primary configuration.}
  \label{tab:ablation}
  \small
  \begin{tabular}{lcccc}
    \toprule
    Config & Hit@10 & MRR & BroadFaith@10 & Description \\
    \midrule
    raw\_kg & 0.60 & 0.35 & 0.00 & No verification \\
    \textbf{+llm} & \textbf{0.70} & \textbf{0.37} & 0.00 & LLM verify, no faith check \\
    \textbf{+llm+faith} & \textbf{0.50} & 0.19 & \textbf{0.98} & \textbf{Primary mode} \\
    +learned & 0.20 & 0.06 & 0.00 & ComplEx score only \\
    full & 0.40 & 0.09 & 0.98 & All modules (incl.\ ComplEx) \\
    no\_kg & 0.43 & 0.13 & 0.00 & KG scaffold off \\
    no\_faith & 0.30 & 0.08 & 0.00 & Module~D off (from full) \\
    \bottomrule
  \end{tabular}
\end{table}

Key findings: (1) \texttt{+llm} achieves the best recall (0.70) but zero faithfulness. Adding Module~D (\texttt{+llm+faith}) reduces recall to 0.50 but enables faithfulness (0.98). This is the core tradeoff: Module~D's faithfulness check does not affect generation (it is post-hoc), but the \texttt{+llm+faith} configuration differs from \texttt{+llm} in that it reports faithfulness labels, which the paper's primary metric requires. (2) ComplEx scoring (\texttt{+learned}) \emph{harms} recall (0.20) because MRR=0.013 means scores are unreliable. (3) \texttt{full} (ComplEx + LLM + faith) underperforms \texttt{+llm+faith} (0.40 vs.\ 0.50), confirming ComplEx is counterproductive on noisy KGs.

\subsection{Faithfulness Validation (Q3)}
\label{sec:exp:faith}

We validate the faithfulness metric using two independent judges:

\begin{table}[tb]
  \centering
  \caption{Faithfulness validation. Two independent judges compared against the original LLM judge.}
  \label{tab:faith_val}
  \small
  \begin{tabular}{lccc}
    \toprule
    Judge pair & BroadFaith@10 & Agreement & Cohen's $\kappa$ \\
    \midrule
    LLM judge (original) & 0.49 & --- & --- \\
    NLI model (DeBERTa-v3-MNLI) & 0.50 & 98\% & 0.00 \\
    LLM proxy (different prompt) & --- & 45\% & 0.095 \\
    \bottomrule
  \end{tabular}
\end{table}

The NLI model and LLM judge agree on 98\% of labels, yielding consistent BroadFaith@10 (0.50 vs.\ 0.49). However, $\kappa{=}0.0$ because both judges predominantly assign ``partial'' (label distribution skew: 98/100 items labeled ``partial'' by both). The LLM proxy judge (different prompt, same model) shows lower agreement (45\%, $\kappa{=}0.095$), suggesting prompt sensitivity. The earlier $\kappa{=}0.095$ (LLM vs.\ proxy) and $\kappa{=}0.0$ (LLM vs.\ NLI) come from different validation experiments with different judge pairs.

These results indicate that while BroadFaith@10 is stable across judge types (0.49--0.50), the \emph{fine-grained labels} (entailed vs.\ partial vs.\ none) are unreliable. This motivates reporting StrictFaith separately (Eq.~\ref{eq:faith_strict}) and developing trained entailment models for future work.

\subsection{KG Cleaning Evaluation (Q4)}
\label{sec:exp:cleaning}

Since we know which triples are synthetic noise (score $<0.5$), we directly evaluate Module~A's cleaning accuracy on 200 sampled triples (100 noise, 100 clean):

\begin{table}[tb]
  \centering
  \caption{KG cleaning evaluation (200 triples: 100 noise + 100 clean). The LLM verifier detects all noise (Recall=1.0) but also flags 95/100 clean triples (high FP rate).}
  \label{tab:cleaning}
  \small
  \begin{tabular}{lcccc}
    \toprule
    Metric & Precision & Recall & F1 & Accuracy \\
    \midrule
    Noise detection & 0.513 & 1.000 & 0.678 & 0.525 \\
    \bottomrule
  \end{tabular}
\end{table}

The LLM verifier achieves perfect recall---it flags all 100 noise triples as unsupported or contradicted. However, precision is only 0.513: 95 of 100 clean triples are also flagged as problematic, yielding 5 true negatives. This high false-positive rate explains why LLM verification helps downstream recall despite imprecise cleaning: the verification step removes far more noise than signal, shifting the signal-to-noise ratio favorably. In the ablation (Table~\ref{tab:ablation}), \texttt{+llm} (verification without faithfulness check) achieves the best Hit@10 (0.70), confirming that aggressive cleaning improves generation quality even when many valid triples are lost.

\subsection{KGQA Generalization}
\label{sec:exp:kgqa}

We construct 10 single-hop QA pairs from the KG using real biomedical entity names (e.g., ``What does \emph{metformin} treat?''). Table~\ref{tab:kgqa} reports exact-match (EM) accuracy. KG-SCoRE and Vanilla LLM both achieve EM=0.10, while ToG and RoG score 0.00. The low scores across all methods reflect the difficulty of exact-string matching against gold answers, but the non-zero results for KG-SCoRE and Vanilla LLM confirm that using real entity names (rather than synthetic IDs) enables the LLM to leverage its parametric knowledge. KG-SCoRE additionally achieves BroadFaith@5 of 1.00, demonstrating that its outputs remain path-consistent even when the answer is incorrect. This confirms the \emph{faithfulness--correctness decoupling}: a method can be perfectly faithful to a KG path while producing a wrong answer.

\begin{table}[tb]
  \centering
  \caption{KGQA with real entity names (10 single-hop questions). EM = exact match.}
  \label{tab:kgqa}
  \small
  \begin{tabular}{lcc}
    \toprule
    Method & EM & BroadFaith@5 \\
    \midrule
    Vanilla LLM & 0.10 & --- \\
    ToG & 0.00 & --- \\
    RoG & 0.00 & --- \\
    \textbf{KG-SCoRE} & 0.10 & \textbf{1.00} \\
    \bottomrule
  \end{tabular}
\end{table}

\subsection{$\tau$ Sensitivity Analysis}
\label{sec:exp:tau}

We sweep the ComplEx confidence threshold $\tau$ over four values to assess its effect on triple filtering and downstream recall (Table~\ref{tab:tau}). At $\tau{\leq}0.25$, no triples are filtered (all ComplEx scores exceed the threshold), yielding identical Hit@10. At $\tau{=}0.40$, 110 low-confidence triples are removed, but Hit@10 does not improve. This confirms that ComplEx scores on noisy KGs (MRR=0.013) do not provide useful signal for triple filtering, reinforcing the ablation finding that learned scoring is counterproductive in this regime.

\begin{table}[tb]
  \centering
  \caption{$\tau$ sensitivity (5 pairs, full pipeline with ComplEx). \#below = triples filtered.}
  \label{tab:tau}
  \small
  \begin{tabular}{lcc}
    \toprule
    $\tau$ & \#below & Hit@10 \\
    \midrule
    0.05 & 0 & 0.20 \\
    0.15 & 0 & 0.20 \\
    0.25 & 0 & 0.00 \\
    0.40 & 110 & 0.20 \\
    \bottomrule
  \end{tabular}
\end{table}

\subsection{Discussion}
The results paint a nuanced picture. Under fair evaluation, KG-SCoRE's faithfulness advantage is real but concentrated in strict faithfulness: ToG and RoG both achieve high BroadFaith (0.94--0.99) and Coverage (1.00), but their StrictFaith is low (0.04--0.05). KG-SCoRE doubles StrictFaith (0.10) while maintaining competitive recall. The ComplEx component is counterproductive on noisy KGs---a key design finding confirmed by both the ablation (Table~\ref{tab:ablation}) and the $\tau$ sensitivity analysis (Table~\ref{tab:tau}): no threshold yields a recall improvement. KG cleaning evaluation reveals that LLM verification achieves perfect noise recall (1.0) but low precision (0.513), explaining why aggressive cleaning shifts the signal-to-noise ratio favorably despite high false positives. The faithfulness metric's fine-grained labels are unreliable ($\kappa{\approx}0$), but the aggregate BroadFaith score is stable across judge types (0.49--0.50). These findings suggest that LLM-based verification is a viable approach for producing traceable hypotheses, but the faithfulness metric requires refinement (stricter labels, trained entailment models).
